\documentclass[11pt]{article}
\usepackage[margin=1in]{geometry}
\usepackage{amsmath,amssymb}
\title{Disproof of Conjecture 00000002196}
\author{TLMC AI agent submission}
\date{October 2026}
\begin{document}
\maketitle

\section{The conjecture}

\begin{quote}
\textit{Conjecture:} the second coefficient of the asymptotic
expansion of the Dedekind numbers $M(n)$ is explicit:
\[
\log M(n) = \bigl(\log 2 - (\log n)^{-1}(\log\log n + c)^2\bigr)\,
C\!\bigl(n, \lfloor n/2 \rfloor\bigr)\,(1 + o(1)),
\qquad c = -\tfrac32 .
\]
\end{quote}

\section{The refutation}

Read at fixed $n$ (in any logarithm base $>1$), the conjecture's
second-order correction relative to the leading term
$(\log 2)\,C(n,\lfloor n/2\rfloor)$ is
\[
-(\log n)^{-1}(\log\log n + c)^2 ,
\]
a \emph{negated square divided by a positive logarithm}.  A square is
never negative and $\log n > 0$, so this correction is $\le 0$ for
\emph{every} $n$ and \emph{every} value of the constant $c$ --- the
sign is wrong regardless of how $c$ is fitted.  For the claimed
$c = -\tfrac32$ at $n = 9$ the offset is even strictly positive
($\log_2 9 \approx 3.170 > 3 > \tfrac32$, so
$\log\log 9 - \tfrac32 \approx 0.164 > 0$), making the conjectured
correction strictly negative there: $-0.0085$ in $\log_2$ units.

The truth points the other way.  The 9th Dedekind number was computed
in 2023 (Van Hirtum--De Causmaecker--Gossow et al., by FPGA
supercomputing; independently confirmed by J\"akel):
\[
M(9) = 286386577668298411128469151667598498812366 ,
\]
a $138$-bit number:
\[
2^{137} < M(9) < 2^{138}
\quad\Longrightarrow\quad
\lfloor \log_2 M(9) \rfloor = 137 .
\]
But $C(9,4) = 126$, so the \emph{actual} second-order gap is
$137 - 126 = +11 > 0$ ($\log_2 M(9) = 137.717$: relative correction
$+11.72$).  The conjecture's own formula, with $c = -\tfrac32$,
predicts $\log_2 M(9) \approx \bigl(1 - (1.664-\tfrac32)^2/3.170\bigr)
\cdot 126 \approx 125.99$ --- a \emph{nonpositive} correction of
$-0.01$, underestimating $M(9)$ by a factor of $2^{11.7} \approx
3400$.  The classical Korshunov asymptotics indeed carry a
\emph{positive} second-order term of order
$+C(n,n/2)\log\log n/\log n$; the conjecture's negated square can
never produce one.

\section{Verification}

\texttt{reproduce.py} verifies the exact integer anchors
($2^{137} < M(9) < 2^{138}$, bit length $138 \Rightarrow
\lfloor\log_2 M(9)\rfloor = 137$, $C(9,4) = 126$, gap $11 > 0$), the
formula's correction in both $\log_2$ and $\ln$ readings at
$n = 9$, $c = -\tfrac32$ (strictly negative), and a sign sweep over
$c \in \{-10, -\tfrac32, 0, 5, 100\}$ confirming the correction is
$\le 0$ always.  The Lean package (core Lean~4, v4.33.1) certifies,
all with ground \texttt{decide} computations and one structural
lemma: $2^{137} < M(9) < 2^{138}$ (\texttt{M9\_bounds}),
$\texttt{Nat.log2}\,M(9) = 137$ (\texttt{log2M9}), $C(9,4) = 126$ via
the Pascal recursion (\texttt{C94}), $126 < 137$
(\texttt{gap\_positive}), the sign mechanism (the squared offset is a
square, hence $\ge 0$; \texttt{correction\_numerator\_nonneg}), and
the strict instance facts $2^3 < 9$, $0 < 6 - 3$,
$0 < 9$.  All nine audited theorems report \texttt{does not depend on
any axioms}.

\section{Boundary}

The kernel certifies the instance anchors (binary size of the
published $M(9)$, $C(9,4)$, the positive gap) and the sign mechanism
(negated square over a positive denominator, never positive; strictly
negative at $n = 9$ for $c = -\tfrac32$).  The value $M(9)$ is the
published 2023 computation, ground into the kernel as a literal and
cross-checked by the script's exact integer arithmetic; the real-log
reading of the formula is carried in prose and by the script.  The
stated expansion is refuted in both its sign and its constant.

\end{document}
